\documentclass{article}

\usepackage{amsmath}
\usepackage{amsthm}
\usepackage{amssymb}

\usepackage[utf8]{inputenc}     
\usepackage[T1]{fontenc}
\usepackage[magyar]{babel}
\usepackage{url}

\title{\textit{\textbf{Szögszámolás}}}
\author{Li Mingdao, Rovenszky Robin}
\date{Legutóbb frissítve: 2026. Május 14.}

\newtheorem{theorem}{Tétel}

\begin{document}

\maketitle
\tableofcontents
\newpage

\begin{theorem}[0. tétel]
    Háromszög szögeinek összege $180^\circ$
\end{theorem}

Az alábbiakban \textbf{irányítatlan szögeket} fogunk használni.
A diszkussziókat pedig nem soroljuk fel itt. \textit{Információ az irányított szögekről: }
\url{https://web.evanchen.cc/handouts/Directed-Angles/Directed-Angles.pdf}
\begin{theorem}[Középponti szögek tétele]
Legyen $A$, $B$, $C$ pontok az $O$ középpontú $\omega$ körön. Tegyük fel, hogy $A$, $O$ a $BC$ egyenes azonos oldalán van. Ekkor $\angle{BOC}=2\angle{BAC}.$
\end{theorem}

\begin{theorem}[Kerületi szögek tétele]
Legyen $A$, $B$, $C$, $D$ pontok az $\omega$ körön. Tegyük fel, hogy $A$, $D$ a $BC$ egyenes azonos oldalán van. Ekkor $\angle{BAC}=\angle{BDC}.$
\end{theorem}

\begin{theorem}[szükséges és elégséges feltétel a 4 pont konciklikusságra]
Legyen $ABCD$ konvex négyszög. $ABCD$ akkor és csak akkor húrnégyszög, ha $\angle{BAD}+\angle{BCD}=180^\circ$.
\end{theorem}

Körív felezőpontjai = Midpoints of arcs\\
The subtended angle of an arc determines the arc's length

\begin{theorem}[Érintő szárú kerületi szögek]
    
\end{theorem}

\end{document}